% GENERATED FILE. Edit canonical lessons, then run npm run generate:books.
% canonical-source-hash: fnv1a64:54985805a682d527
% canonical-lessons: LA-C154-iungo, LA-C154-visito, LA-C154-canto, LA-C154-pingo, LA-C154-duco

\chapter{To join, To visit, To sing, To paint, To lead}
\label{ch:la-to-join-to-visit-to-sing-to-paint-to-lead}
\hlchaptermodality{154}

\begin{chapteropening}
\textbf{By the end of this chapter:} \emph{I can say to join, to visit, to sing, to paint and to lead.}

Complete the last lesson of chapter 154: I can say to join, to visit, to sing, to paint and to lead.
\end{chapteropening}

\section[iungō, iungere]{iungō, iungere — to join}

\label{lesson:LA-C154-iungo}



\begin{quote}
Before the new one: say the Latin for to climb, then the Latin for to go down.
\end{quote}

\subsection*{You'll want to know: iungō, iungere}
\textbf{iungō, iungere} — \textquotedblleft{}to join\textquotedblright{}.

\begin{cousinweb}
\textbf{iungō} is inherited, usually traced to Proto-Indo-European \emph{*yewg-}, \textquotedblleft{}to join, to yoke\textquotedblright{} (EtymDB 2.1); the n inside the present stem is an old mark of the verb. Greek \emph{zeugnymi}, \textquotedblleft{}to yoke\textquotedblright{}, and English \emph{yoke} are cousins (Lewis \& Short).
\end{cousinweb}

\begin{grammarlens}[title={What you've built}]
One more word for this chapter.
\end{grammarlens}

\subsection*{Your turn}
\begin{itemize}
\raggedright
  \item \emph{Say it:} \emph{iungō}
  \item \emph{Say it:} \emph{iungō}, once more
  \item \emph{Say it:} say \emph{iungō}
  \item \emph{Recall:} say the Latin for to climb, then the Latin for to go down, then say \emph{iungō} again
\end{itemize}

\begin{tcolorbox}[breakable,colback=teal!4,colframe=teal!35!black,title={Before you move on}]
What does iungō, iungere mean? (\textquotedblleft{}To join\textquotedblright{}.) Say it once more.
\end{tcolorbox}
\section[vīsitō, vīsitāre]{vīsitō, vīsitāre — to visit}

\label{lesson:LA-C154-visito}



\begin{quote}
Before the new one: say the Latin for to go down, then the Latin for to join.
\end{quote}

\subsection*{You'll want to know: vīsitō, vīsitāre}
\textbf{vīsitō, vīsitāre} — \textquotedblleft{}to visit\textquotedblright{}.

\begin{cousinweb}
\textbf{vīsitō} is a frequentative of \emph{vīsō}, \textquotedblleft{}to go and see\textquotedblright{} (Lewis \& Short), itself built on \emph{videō}, \textquotedblleft{}to see\textquotedblright{}: to keep going to see.
\end{cousinweb}

\begin{grammarlens}[title={What you've built}]
One more word for this chapter.
\end{grammarlens}

\subsection*{Your turn}
\begin{itemize}
\raggedright
  \item \emph{Say it:} \emph{vīsitō}
  \item \emph{Say it:} \emph{vīsitō}, once more
  \item \emph{Say it:} say \emph{vīsitō}
  \item \emph{Recall:} say the Latin for to go down, then the Latin for to join, then say \emph{vīsitō} again
\end{itemize}

\begin{tcolorbox}[breakable,colback=teal!4,colframe=teal!35!black,title={Before you move on}]
What does vīsitō, vīsitāre mean? (\textquotedblleft{}To visit\textquotedblright{}.) Say it once more.
\end{tcolorbox}
\section[cantō, cantāre]{cantō, cantāre — to sing}

\label{lesson:LA-C154-canto}



\begin{quote}
Before the new one: say the Latin for to join, then the Latin for to visit.
\end{quote}

\subsection*{You'll want to know: cantō, cantāre}
\textbf{cantō, cantāre} — \textquotedblleft{}to sing\textquotedblright{}.

\begin{cousinweb}
\textbf{cantō} is a frequentative of \emph{canō}, \textquotedblleft{}to sing\textquotedblright{} (Lewis \& Short), built on its stem \emph{cant-}, and in everyday Latin it came to replace \emph{canō}. \emph{canō} is usually traced to Proto-Indo-European \emph{*keh2n-}, \textquotedblleft{}to sing\textquotedblright{} (EtymDB 2.1).
\end{cousinweb}

\begin{grammarlens}[title={What you've built}]
One more word for this chapter.
\end{grammarlens}

\subsection*{Your turn}
\begin{itemize}
\raggedright
  \item \emph{Say it:} \emph{cantō}
  \item \emph{Say it:} \emph{cantō}, once more
  \item \emph{Say it:} say \emph{cantō}
  \item \emph{Recall:} say the Latin for to join, then the Latin for to visit, then say \emph{cantō} again
\end{itemize}

\begin{tcolorbox}[breakable,colback=teal!4,colframe=teal!35!black,title={Before you move on}]
What does cantō, cantāre mean? (\textquotedblleft{}To sing\textquotedblright{}.) Say it once more.
\end{tcolorbox}
\section[pingō, pingere]{pingō, pingere — to paint}

\label{lesson:LA-C154-pingo}



\begin{quote}
Before the new one: say the Latin for to visit, then the Latin for to sing.
\end{quote}

\subsection*{You'll want to know: pingō, pingere}
\textbf{pingō, pingere} — \textquotedblleft{}to paint\textquotedblright{}.

\begin{cousinweb}
\textbf{pingō} is inherited, usually traced to Proto-Indo-European \emph{*peyk-}, \textquotedblleft{}to mark, to colour\textquotedblright{} (EtymDB 2.1); the n inside the present stem is missing from its participle \emph{pictus}, \textquotedblleft{}painted\textquotedblright{}.
\end{cousinweb}

\begin{grammarlens}[title={What you've built}]
One more word for this chapter.
\end{grammarlens}

\subsection*{Your turn}
\begin{itemize}
\raggedright
  \item \emph{Say it:} \emph{pingō}
  \item \emph{Say it:} \emph{pingō}, once more
  \item \emph{Say it:} say \emph{pingō}
  \item \emph{Recall:} say the Latin for to visit, then the Latin for to sing, then say \emph{pingō} again
\end{itemize}

\begin{tcolorbox}[breakable,colback=teal!4,colframe=teal!35!black,title={Before you move on}]
What does pingō, pingere mean? (\textquotedblleft{}To paint\textquotedblright{}.) Say it once more.
\end{tcolorbox}
\section[dūcō, dūcere]{dūcō, dūcere — to lead}

\label{lesson:LA-C154-duco}



\begin{quote}
Before the new one: say the Latin for to sing, then the Latin for to paint.
\end{quote}

\subsection*{You'll want to know: dūcō, dūcere}
\textbf{dūcō, dūcere} — \textquotedblleft{}to lead\textquotedblright{}.

\begin{cousinweb}
\textbf{dūcō} is inherited, usually traced to Proto-Indo-European \emph{*dewk-}, \textquotedblleft{}to lead, to pull\textquotedblright{} (EtymDB 2.1); Gothic \emph{tiuhan}, \textquotedblleft{}to lead, to draw\textquotedblright{}, is a cousin (Lewis \& Short), and so are English \emph{tug} and \emph{tow}.
\end{cousinweb}

\begin{grammarlens}[title={What you've built}]
One more word for this chapter.
\end{grammarlens}

\subsection*{Your turn}
\begin{itemize}
\raggedright
  \item \emph{Say it:} \emph{dūcō}
  \item \emph{Say it:} \emph{dūcō}, once more
  \item \emph{Say it:} say \emph{dūcō}
  \item \emph{Recall:} say the Latin for to sing, then the Latin for to paint, then say \emph{dūcō} again
\end{itemize}

\begin{tcolorbox}[breakable,colback=teal!4,colframe=teal!35!black,title={Before you move on}]
What does dūcō, dūcere mean? (\textquotedblleft{}To lead\textquotedblright{}.) Say it once more.
\end{tcolorbox}
